\section{Blocked samplers and the association parameter}\label{sec:blocked}

The location degeneracy does not involve $\alpha$. Yet on $\alpha$,
\texttt{jmjax}'s NUTS was 3--8 times more efficient per second than
JMbayes2 even without the rotation: 5.7--8.1 times across a replicated
simulation sweep with 1 to 12 covariates, and 3.2 and 7.6 times on the two
datasets of Section~\ref{sec:realdata}. JMbayes2 updates the parameters in
blocks, with adaptive random-walk Metropolis steps for the coefficients.
This section shows that most of the difference is a property of blocked
updating in joint models rather than of either implementation.

Write $\psi$ for every parameter other than $\alpha$: the random effects,
$\beta$, $\gamma$, the baseline-hazard coefficients and the variance
components.

\begin{proposition}[Blocked updates of $\alpha$]\label{prop:blocked}
Consider the two-block sampler that draws $\alpha$ exactly from
$p(\alpha \mid \psi, y)$ and then $\psi$ exactly from $p(\psi \mid
\alpha, y)$, and let
\[
r = \frac{\Var\{\E(\alpha \mid \psi, y)\}}{\Var(\alpha \mid y)}
= 1 - \frac{\E\{\Var(\alpha \mid \psi, y)\}}{\Var(\alpha \mid y)},
\]
the fraction of $\alpha$'s posterior variance explained by the other
parameters, and $R^2 \leq r$ the squared multiple correlation of $\alpha$
on $\psi$.
\begin{enumerate}
\item[(i)] At stationarity the lag-1 autocorrelation of the $\alpha$ chain
is $r$.
\item[(ii)] Its integrated autocorrelation time is at least $(1 + r)/(1 -
r)$, so its effective sample size per draw is at most $(1 - r)/(1 + r)
\leq (1 - R^2)/(1 + R^2)$.
\item[(iii)] If the posterior is Gaussian, $r = R^2$, the $\alpha$ chain
is AR(1) with coefficient $R^2$, and the bound in (ii) is attained.
\end{enumerate}
\end{proposition}

\begin{proof}
The $\alpha$ chain of a two-block sampler is a data-augmentation chain.
Its lag-1 autocovariance is $\Var\{\E(\alpha \mid \psi, y)\}$ and its
transition operator is self-adjoint and positive on $L^2$ of the
posterior of $\alpha$ \citep{liu1994}, which gives (i); the second
expression for $r$ is the law of total variance, and $R^2 \leq r$ because
the best linear predictor of $\alpha$ explains no more variance than
$\E(\alpha \mid \psi, y)$. For (ii), the spectral theorem writes the
lag-$k$ autocorrelation of $\alpha$ as $\rho_k = \int \lambda^k\,
\mu(d\lambda)$ for a probability measure $\mu$ on $[0, 1]$, so Jensen's
inequality gives $\rho_k \geq \rho_1^k = r^k$ and $1 + 2\sum_{k \geq 1}
\rho_k \geq (1 + r)/(1 - r)$; the function $(1 - x)/(1 + x)$ is
decreasing. For (iii), $\E(\alpha \mid \psi)$ is linear in $\psi$ and
$\E(\psi \mid \alpha)$ linear in $\alpha$; composing the two regressions
gives $\alpha_{t+1} = c + R^2 \alpha_t + \varepsilon_t$ with
$\varepsilon_t$ independent of the past, since the coefficient is
$\operatorname{Cov}(\alpha, \psi)\,\Sigma_\psi^{-1}
\operatorname{Cov}(\psi, \alpha) / \Var(\alpha) = R^2$, and then
$\rho_k = (R^2)^k$.
\end{proof}

Part (ii) is what makes a measured linear $R^2$ usable: it bounds the
efficiency of any exact two-block sampler from above, whatever the shape
of the posterior. The proposition is an application of known theory; the
two-block sampler is a benchmark, not a model of JMbayes2. Finer blocking
and random-walk Metropolis steps can be expected to make a real
implementation slower still \citep{liu1994,roberts1997}, though neither is
guaranteed to in every case. What the proposition makes precise is the
contrast with HMC, which moves $\alpha$ and $\psi$ together. Under a
fixed metric, a correlation between them is part of the posterior's shape
and is paid for in leapfrog steps, which grow with the square root of the
condition number \citep{neal2011}. A single correlation at level $R^2$
contributes a condition number $(1 + \sqrt{R^2})/(1 - \sqrt{R^2})$, so the
cost grows at most like $1/\sqrt{1 - R^2}$, against $1/(1 - R^2)$ for the
blocked sampler.
Often it grows not at all, because in a joint model the step size is
already set by some other narrow direction, such as the spline
coefficients or the variance components. In a Gaussian check, comparing a
simulated two-block Gibbs sampler with a stylised model of HMC (exact
Hamiltonian dynamics, random integration time), the two-block sampler's
lag-1 autocorrelation
matched $R^2$ to within 0.004, and its integrated autocorrelation time
went 3.0, 9.0, 27.6, 39.0 and 199 as $R^2$ went 0.5, 0.8, 0.93, 0.95 and
0.99. The HMC model's stayed between 1.6 and 2.1, at a cost of 4.8 to 46.8
gradient evaluations per independent draw; with an unrelated narrow ridge
already present, the step count was 22.2 at every $R^2$ from 0 to 0.99.

\subsection{How large $R^2$ is in a joint model}

Two things make $R^2$ large. The first is errors in variables: given the
random effects, each marker trajectory $m_i(t)$ is known and $\alpha$ is
the coefficient of a known time-varying covariate, which the data pin
sharply; marginally, $\alpha$ also carries the uncertainty in every
trajectory. The second is the level of the hazard. Shifting $m_i(t)$ by a
constant for every subject changes $\alpha\, m_i(t)$ by an amount the log
baseline hazard absorbs, so $\alpha$ is correlated with the baseline-hazard
level unless the marker is centred.

We estimated $R^2$ by regressing $\alpha$ on every other sampled quantity
in a long rotated fit (4 chains of 5000 draws; 960 regressors for
\texttt{aids} and 649 for \texttt{pbc2}, adjusted for their number), and
split it by source. On \texttt{aids}, $R^2 = 0.873$; the survival block
alone ($\gamma$ and the baseline-hazard coefficients) explained 0.86 of
$\alpha$'s posterior variance and the random effects alone 0.30. On
\texttt{pbc2} the figures were 0.853, 0.82 and 0.20. The hazard level, not
the random effects, is what a blocked sampler mostly conditions away.
By Proposition~\ref{prop:blocked}(ii), the linear $R^2$ is a lower bound
on the fraction $r$, so the efficiency it implies for an ideal blocked
sampler is an upper bound.

\begin{figure}[t]
\centering
\includegraphics[width=0.66\textwidth]{fig3_alpha.pdf}
\caption{Effective sample size per draw for $\alpha$. Curve: an ideal
two-block sampler, $(1 - R^2)/(1 + R^2)$ (Proposition~\ref{prop:blocked}).
Black open markers: the ideal sampler at the measured $R^2$ (0.873 for
\texttt{aids}, 0.853 for \texttt{pbc2}). JMbayes2 does two to three times
worse than that ideal; the NUTS draws are effectively independent. Green
open markers: the ideal if $\alpha$ were updated in one block with the
survival parameters, so that only the random effects are conditioned away
($R^2 = 0.30$ and $0.20$).}
\label{fig:alpha}
\end{figure}

Figure~\ref{fig:alpha} puts the measurements against the prediction. At
the measured $R^2$ an ideal two-block sampler keeps at most 0.068
(\texttt{aids}) and 0.079 (\texttt{pbc2}) effective draws per draw, one
in 13--15.
JMbayes2 achieved 0.035 and 0.028, and \texttt{jmjax}'s NUTS 1.05 and 1.72
in the same long fit. On a log scale, the ideal benchmark accounts for
80\% of the measured ESS-per-draw gap on \texttt{aids} (15.5 of 30 times)
and 75\% on \texttt{pbc2} (21.7 of 61 times); the rest, a factor of two to
three, is attributable to the implementation. A blocked sampler could
escape most of the structural part by updating $\alpha$ jointly with
$\gamma$ and the baseline-hazard coefficients: with only the random
effects conditioned away, the ideal ESS per draw would be about 0.54 and
0.67. But that block has 10--20 strongly correlated dimensions without
exact conditional draws, and an adaptive random-walk proposal on it is
itself slow. NUTS obtains the joint move without any block design.
